\section{Source version and scope of the integration}
The previous ES SplitZero note retained supported zero, the original incoming
relations, and the exact mass of every bounded divisor state. We now integrate
three later source interfaces: support-changing quotients and both transition
defects (PR28, equations I1--I9); canonical metric secants with their original
relation norms (PR29, sections 3--4); and the new boundary-kernel comparison
(PR30, equations BR1--BR7). The first two are read at Zeta main
\texttt{35b7aeaff0b63b0db71f26c5b921cdbbec2507fb}; the third is read at
\texttt{f57532b97c5c92a8c197e6cd4fef2b36b2d33e97}, an unmerged draft at intake
\cite{SZ28,SZ29,SZ30}. These are not all assigned to the frozen R62 publication.
The 821-page edition is not represented as fully audited. Its periodized theta
integrals and growing-family estimates are not identified with the finite norms
below. The supplied source records formal checks; no new Lean build is claimed.

The new arithmetic constructions are a prime-labelled boundary module, its
socle-to-augmentation comparison, and two canonical moment bounds computed
without selecting successful atoms first. All proofs below are elementary and
include singular finite moment forms. The positive-definite inverse formulas of
the source cannot simply be used on a singular matrix.

\section{Complete original arithmetic, including failed states}
Fix a prime $p\equiv1\pmod4$, an integer $p/4<a<p/2$, and $R=4a-p$.
Write $a=\prod_i q_i^{e_i}$ with distinct increasing primes. A canonical labelled
state is $\alpha=(p,a,R,c,\beta)$, where $c\in\{E,M\}$ and
$\beta\in\prod_i[-e_i,e_i]_{\mathbb Z}$. Set
\[
 u_\alpha=\prod_iq_i^{e_i+\beta_i},\qquad
 g_E=4u+1,\quad g_M=u+a,\qquad D_\alpha=R/\gcd(R,g_c).
 \tag{A1}
\]
All actual prime-factor multiplicities and both channels are retained. For a
finite selection of shells use their disjoint union, not a deduplicated residue
set. The full first-half interval is sufficient by the existing workbench
smallest-denominator theorem; none of the new forcing lemmas requires that
completeness theorem for its validity on a specified finite selection.

\begin{lemma}[Rational candidate and exact return]
Each state gives the positive rational triple
\[
 Q_E=\left(a,\frac{pa+a^2/u}{R},\frac{pa+p^2u}{R}\right),\qquad
 Q_M=\left(a,\frac{p(a+a^2/u)}R,\frac{p(a+u)}R\right).
 \tag{A2}
\]
It solves $4/p=\sum_iQ_i^{-1}$. The last two denominators in reduced form are
both exactly $D_\alpha$. Consequently the candidate is integral precisely when
$D_\alpha=1$; otherwise $D_\alpha\ge3$ is odd. With $x_\alpha=1/D_\alpha$,
\[
 x_\alpha\in(0,1/3]\cup\{1\}.
 \tag{A3}
\]
\end{lemma}
\begin{proof}
We have $\gcd(pa,R)=1$ and $u\mid a^2$. Direct multiplication gives
$(RY-pa)(RZ-pa)=p^2a^2$ for each displayed triple. Expanding gives the
reciprocal identity. Modulo $R$, the two exterior numerators, after multiplying
the first by the unit $u$, are $a^2(4u+1)$ and $4a^2(4u+1)$. The two middle
numerators similarly have gcd with $R$ equal to $\gcd(R,u+a)$. This proves the
exact reduced denominators. Oddness follows from $R\equiv3\pmod4$.
\end{proof}

These are the inherited Type I/II divisor coordinates \cite{ET}. At a hit let
$d_0=\gcd(a,u)$ and retain
\[
 (h,r,s)=(d_0^2/u,u/d_0,a/d_0),\quad a=hrs,\quad u=hr^2,
\]
\[
 E:\quad\kappa=(pr+s)/R,\ (a,hs\kappa,phr\kappa),\qquad
 M:\quad\lambda=(r+s)/R,\ (a,phs\lambda,phr\lambda).
 \tag{A4}
\]
Prime valuations prove integrality and $\gcd(r,s)=1$; clearing reciprocals
verifies (A4). The ordered inverse is
\[
 u=\frac{a^2}{RY-pa}\ (E),\qquad
 u=\frac{pa^2}{RY-pa}\ (M).
 \tag{A5}
\]
The exterior ordering bit, if both raw orders are wanted, and the middle
involution $u\mapsto a^2/u$ are not identified. Our counts use the canonical
$E,M,u$ labels just declared, not distinct unordered denominator triples.

\section{Prime-labelled boundary modules and the success comparison}
Let $\mathcal L$ be the finite set of primes dividing the selected shell
moduli. For every state define
\[
 d_{\alpha\ell}=\max\{v_\ell(R_\alpha)-v_\ell(g_\alpha),0\},\qquad
 D_\alpha=\prod_{\ell\in\mathcal L}\ell^{d_{\alpha\ell}}.
 \tag{B1}
\]
A prime absent from a shell has exponent zero. Over
$A=\mathbb Q[z_\ell:\ell\in\mathcal L]$, put
\[
 J_\alpha=A/(z_\ell^{d_{\alpha\ell}+1}:\ell\in\mathcal L),\quad
 \mathcal J=\bigoplus_\alpha J_\alpha,\quad
 V=\bigoplus_\alpha\mathbb Q e_\alpha.
 \tag{B2}
\]
These are independent variables for different primes. The construction is
not a tensor of cones with an unspecified common parameter.

\begin{theorem}[The retained boundary comparison]
The joint annihilator $\operatorname{Soc}\mathcal J=\bigcap_\ell\ker z_\ell$
has the labelled basis $z^{d_\alpha}e_\alpha$. The maps
\[
 s:V\xrightarrow{\sim}\operatorname{Soc}\mathcal J,
 \quad e_\alpha\mapsto z^{d_\alpha}e_\alpha,\qquad
 \epsilon:\mathcal J\to V,\quad F_\alpha\mapsto F_\alpha(0)e_\alpha
 \tag{B3}
\]
obey
\[
 \boxed{\epsilon s=\operatorname{diag}(\mathbf1_{D_\alpha=1}).}
 \tag{B4}
\]
In particular the rank and trace of this comparison, and the dimension of the
multidegree-zero socle, equal the original success count $H$.
\end{theorem}
\begin{proof}
The residue monomials $z^b$, $0\le b_\ell\le d_{\alpha\ell}$, form a basis.
Multiplication by $z_\ell$ sends distinct surviving basis vectors to distinct
basis vectors, so its kernel consists exactly of those with
$b_\ell=d_{\alpha\ell}$. Intersecting leaves the single top monomial.
Its coefficient supplies the inverse to $s$. Evaluation at zero sends this
monomial to one if all exponents vanish and to zero otherwise, proving (B4).
\end{proof}

The direct source of (B3) is BR1 in \cite{SZ30}: if multiplication by $v$ and
$f$ are injective on $M$, then
$M/vM\simeq\ker(v:M/vf(M)\to M/vf(M))$ via $[x]\mapsto[f(x)]$.
Well-definedness uses linearity, injectivity uses cancellation by $f$ before
quotienting, and surjectivity uses cancellation by $v$ on $M$.
Here each single-prime instance has $M=\mathbb Q[z_\ell]$,
$v=z_\ell$, and $f=z_\ell^{d_{\alpha\ell}}$; the joint statement was proved
above without an unproved tensor-homology identification. Classical Koszul
regularity is compatible with this polynomial model \cite{Stacks}; no general
Tor functor is needed for the displayed computation.

Every labelled summand of the full socle is present, also at failures. Its
ordinary augmentation and its graded kernel are different observations.
In the original SplitZero carrier, a zero image has the same nonbottom support
label: it is $e$ in that fibre, not external absence $\tau$. This instantiates
the receiving-zero criterion I2 and the boundary-image criterion BR7
\cite{SZ28,SZ30}. The joint annihilator is stated explicitly; it is not silently
the annihilator of one arbitrary scalar.

\paragraph{Reduction and its opposite-direction map.}
For $d'\le d$ coordinatewise, reduction $Q:J_d\to J_{d'}$ is the usual
monomial truncation. Its kernel is the span of monomials exceeding at least
one $d'_\ell$. It sends the top monomial to zero unless $d=d'$.
In the opposite direction,
\[
 I:J_{d'}\hookrightarrow J_d,\qquad F\mapsto z^{d-d'}F
 \tag{B5}
\]
is injective and preserves the top socle; its inverse on its image subtracts
the exponent vector $d-d'$. The composite $QI$ is multiplication by
$z^{d-d'}$ modulo the smaller ideal, not generally the identity.
Reduction from a modulus $R$ to a divisor of $R$ decreases the vector (B1),
so it uses exactly $Q$. This is a local observation of the same marked
arithmetic leaf, not a claim that the reduced modulus defines a new ES shell.
Independent local successes cannot be glued by
replacing this map by the socle-preserving reverse map. Arithmetic coefficient
multiplication by $x_d=\prod\ell^{-d_\ell}$ has transition defect
\[
 Q M_{x_d}-M_{x_{d'}}Q=(x_d-x_{d'})Q.
 \tag{B6}
\]
This records the extra coefficient defect of I9 in \cite{SZ28}; it is not
assumed zero.

\section{The original split weights commute with this boundary model}
Split $e=f+b$. The map $q:(\xi,\eta,c)\mapsto(\xi+\eta,c)$ has complete fibre
\[
 \max(-f_i,\beta_i-b_i)\le\xi_i\le\min(f_i,\beta_i+b_i),\quad
 \eta=\beta-\xi,
\]
and cardinality
\[
 n(\beta)=\prod_i\bigl(\min(f_i,\beta_i+b_i)-\max(-f_i,\beta_i-b_i)+1\bigr).
 \tag{W1}
\]
Every interval is nonempty because $|\beta_i|\le f_i+b_i$.
The two canonical channel labels remain separate even if numerical values
coincide. At every split occurrence use the same $d_\alpha,J_\alpha,s,\epsilon$
as its returned state. Pullback $U$ duplicates coefficients, and
\[
 (EF)_\alpha=\frac1{n(\beta)}\sum_{q(\omega)=\alpha}F_\omega.
 \tag{W2}
\]
Then $EU=1$, both maps commute with all $z_\ell$, with $s,\epsilon$, and with
$x_\alpha$. The kernel of $E$ consists of fibre sums zero; its complete
primitive is the nonreference-coordinate star-difference expansion. Native
occurrence counting is not the original arithmetic metric.

For original value-vector weights $w_\alpha>0$, use occurrence weights
$w_\alpha/n(\beta)$. The maps (W2) satisfy $U^*=E$ and $U^*U=1$ and
\[
 \|F\|^2=\|EF\|^2+\|F-UEF\|^2.
 \tag{W3}
\]
The proof groups the sums by fibres; the cross term vanishes because each
remainder fibre sum is zero. On coefficient vectors with pushforward instead
of pullback the dual convention is cost $n(\beta)$ and lift $1/n(\beta)$;
these types are related by the diagonal fibre-size map, not conflated.
For every polynomial $P$, in particular all moment and Gram entries,
\[
 \sum_\alpha w_\alpha P(x_\alpha)
 =\sum_\omega\frac{w_{q\omega}}{n(\beta)}P(x_{q\omega}).
 \tag{W4}
\]
Arbitrary further source transports retain their off-diagonal pairings, as
required by BR8 and SP1--SP13 \cite{SZ30,SP}. Different labels alone do not
prove orthogonality after a map identifies them.

\paragraph{The original residue cohomology and the arithmetic image are distinct.}
For one actual split and channel, let $\mathscr P$ be the matching-pair set,
$S$ the original successful exponent labels, and $I$ the common residue set.
For the two actual maps $X\to Z\leftarrow Y$, use the complex
$B_X\oplus B_Y\hookrightarrow\mathbb Q^X\oplus\mathbb Q^Y\to\mathbb Q^Z$,
where $B_X,B_Y$ are generated by the nonreference star differences in each
residue fibre and the last map is left fibre sum minus right fibre sum.
A cycle has equal fibre sums; subtracting the two reference representatives
leaves the specified star relations. Thus its complete degree-zero quotient
is $\mathbb Q^I$, including its inverse and all boundary primitives.
This is the preceding residue cohomology. In contrast, (B4)
has image $\mathbb Q^S$. Their exact comparison is the span
\[
 \mathbb Q^S\xleftarrow{q_*}\mathbb Q^{\mathscr P}
 \xrightarrow{r_*}\mathbb Q^I.
 \tag{W5}
\]
Both maps sum coefficients over their indicated fibres. Their kernels have
the complete nonreference star-difference bases. A map induced directly from
$q_*$ to $r_*$ exists only when the residue is constant on every arithmetic
fibre: this is equivalent to $\ker q_*\subseteq\ker r_*$. It is not assumed.
The retained arithmetic section $\sigma_q e_\alpha=
\sum_{q\omega=\alpha}e_\omega/n(\beta)$ defines an actual comparison
\[
 B=r_*\sigma_q,\qquad B_{z\alpha}=
 \frac{|r^{-1}(z)\cap q^{-1}(\alpha)|}{n(\beta)},\qquad
 \mathbf1^{\mathsf T}B=\mathbf1^{\mathsf T}.
 \tag{W6}
\]
Thus it preserves total arithmetic mass, not necessarily vector dimension
or the original metric. For the previous counting-source harmonic metric
$G_I=\operatorname{diag}(1/|X_z|+1/|Y_z|)$, the transported form is exactly
$B^{\mathsf T}G_IB$, including all cross terms.

At $p=2521,a=636,R=23$, $f=(1,1,1),b=(1,0,0)$, there are five matched
pairs, four common channel/residues, and three original channel/divisor
states. With columns $(E,477),(M,8),(M,50562)$ and rows
$(E,13),(M,11),(M,21),(M,22)$,
\[
 B=\begin{pmatrix}1&0&0\\0&1/2&0\\0&0&1/2\\0&1/2&1/2\end{pmatrix},
 \quad G_I=\operatorname{diag}(2,2,2,3/2),\quad
 B^{\mathsf T}G_IB=
 \begin{pmatrix}2&0&0\\0&7/8&3/8\\0&3/8&7/8\end{pmatrix}.
 \tag{W7}
\]
The mixed $3/8$ entries are retained. Consequently the new success-projector
trace is three, not an identification with the old residue cohomology's
dimension four. Both are quotients of the same fully marked pair space,
with their own kernel, primitive, and norm.

\section{Canonical two-sided bounds from the unchanged moment data}
The following result applies to any finite positive measure supported on
$(0,c]\cup\{1\}$, where $0<c<1$. In the ES instance use $c=1/3$ and unit
weight on every original state. Let
\[
 m_j=\sum_\alpha w_\alpha x_\alpha^j,\quad
 W=m_0,\quad H=\sum_{x_\alpha=1}w_\alpha.
\]
Do not divide by $W$: the original mass is retained. For degree $d\ge0$ set
\[
 U_d=\min_{\deg P\le d,\ P(1)=1}\sum_\alpha w_\alpha P(x_\alpha)^2,
\quad
 L_d=\max_{\deg P\le d,\ P(1)=1}
 \frac1{1-c}\sum_\alpha w_\alpha(x_\alpha-c)P(x_\alpha)^2.
 \tag{M1}
\]
The first is the classical Christoffel variational problem \cite[eqs.~(2.2)--(2.3)]{Las}.
We retain its singular finite-source cases and the signed lower counterpart.
These are not an identification with the source's analytic theta moment form.

Write $P=1+(x-1)q(x)$ with $q$ of degree at most $d-1$.
For $0\le i,j<d$ define
\[
 (G_+)_{ij}=m_{i+j+2}-2m_{i+j+1}+m_{i+j},\quad
 (b_+)_i=m_{i+1}-m_i,
\]
\[
 (G_-)_{ij}=-m_{i+j+3}+(c+2)m_{i+j+2}
 -(2c+1)m_{i+j+1}+cm_{i+j},
\quad
 (b_-)_i=m_{i+2}-(c+1)m_{i+1}+cm_i.
 \tag{M2}
\]

\begin{theorem}[Canonical bracket with complete singular fibres]
Both $G_+,G_-$ are positive semidefinite and $b_\pm$ belong to their ranges.
Choose any exact solutions $G_\pm y_\pm=b_\pm$. Then
\[
 \boxed{U_d=m_0-b_+^{\mathsf T}y_+,\qquad
 L_d=\frac{m_1-cm_0+b_-^{\mathsf T}y_-}{1-c}.}
 \tag{M3}
\]
All optimizing polynomials are
\[
 P_+=1-(x-1)y_+(x)+(x-1)\ker G_+,\qquad
 P_-=1+(x-1)y_-(x)+(x-1)\ker G_-.
 \tag{M4}
\]
They obey
\[
 \boxed{L_d\le H\le U_d,\qquad L_d\nearrow H,\quad U_d\searrow H.}
 \tag{M5}
\]
For a finite measure equality holds at some finite degree. When original
weights count states, $H$ is an integer, and $\lceil L_d\rceil=\lfloor U_d\rfloor$
is an exact count certificate. In particular $L_d>0$ forces a witness and
$U_d<1$ certifies no witness in the stated finite domain.
\end{theorem}
\begin{proof}
The two Grams are sums of rank-one matrices with coefficients
$w(x-1)^2$ and $w(c-x)(x-1)^2$, respectively. The second coefficient is
nonnegative on the declared support, including its zero value at $x=1$.
A kernel polynomial of either Gram vanishes at every point carrying a
positive corresponding coefficient. The same points carry its $b$ terms;
at $x=1$, and additionally at $x=c$ for the second form, those terms vanish.
Thus each $b$ annihilates the kernel, proving membership in the range.
Expanding (M1) in $q$ gives the upper expression
$m_0+2b_+^{\mathsf T}q+q^{\mathsf T}G_+q$ and the lower numerator
$m_1-cm_0+2b_-^{\mathsf T}q-q^{\mathsf T}G_-q$.
Completing the squares proves (M3)--(M4), including attainment and all fibres.
Successful atoms contribute exactly their original mass to either expression.
Every other contribution is nonnegative for the upper and nonpositive for
the lower. This proves the bounds. Increasing degree enlarges the feasible
polynomial space. A polynomial equal to one at 1 and vanishing at every other
support point gives equality, proving finite termination.
\end{proof}

\paragraph{The relation norms, not only their zero quotient.}
For compatible choices of optimizers at $i\le j$, the original secants are
\[
 U_i-U_j=\sum_\alpha w_\alpha(P_{+,i}-P_{+,j})^2(x_\alpha),
\]
\[
 L_j-L_i=\frac1{1-c}\sum_\alpha w_\alpha(c-x_\alpha)
              (P_{-,i}-P_{-,j})^2(x_\alpha).
 \tag{M6}
\]
Both follow by orthogonality of the optimizer to all allowed corrections.
The difference has explicit primitive $(P_i-P_j)/(x-1)$.
The second right side is nonnegative because the difference vanishes at 1.
At $x=c$ its observation weight is zero: those original states remain present
in the full source and in the upper form. They are not declared external zero
or silently identified with the source's original relation module. This is the
finite arithmetic implementation of the metric-boundary discipline in
\cite[sections 3--4]{SZ29}.

\paragraph{The first-degree test as four integer gcd moments.}
At one shell, set $S_j=\sum_\alpha\gcd(R,g_\alpha)^j$ for $0\le j\le3$;
then $m_j=S_j/R^j$. Put
\[
 A_0=3S_1-RS_0,\quad A_1=3S_2-RS_1,\quad A_2=3S_3-RS_2.
 \tag{M8}
\]
The degree-one signed localizing matrix is
\[
 \begin{pmatrix}m_1-m_0/3&m_2-m_1/3\\
 m_2-m_1/3&m_3-m_2/3\end{pmatrix}.
\]
It is negative semidefinite exactly when
$A_0\le0$, $A_2\le0$, and $A_0A_2\ge A_1^2$.
A positive direction cannot be a polynomial vanishing at one, since its
remaining terms would all be nonpositive. It can therefore be normalized
and is one of the competitors in (M1). Conversely $L_1>0$ supplies that
direction. We have proved the exact criterion
\[
 \boxed{L_1>0\iff A_0>0\ \text{or}\ A_2>0\ \text{or}\ A_0A_2<A_1^2.}
 \tag{M9}
\]
Every empty shell must satisfy the three opposite inequalities. For example, the measure $\delta_1+10^4\delta_{1/5}+10^4\delta_{1/7}$ has
$H=1$ but $L_1=-8189/1811<0$. Thus the general
finite-support theorem allows occupied measures with $L_1\le0$;
no all-prime sufficiency of degree one is inferred from the finite ES scan.

\paragraph{Uniform finite stopping bound.}
For $c=1/3$, use the existing comparison polynomial
$C_d(x)=T_d(6x-1)/T_d(5)$, with the usual Chebyshev recurrence. Its value at one
is one and its absolute value on $[0,1/3]$ is at most $T_d(5)^{-1}$.
Using it as a competitor proves
\[
 0\le U_d-H\le\frac{W}{T_d(5)^2},\qquad
 0\le H-L_d\le\frac{W}{2T_d(5)^2}.
 \tag{M7}
\]
Hence $T_d(5)^2>3W/2$ guarantees a unique integer in the bracket. Since
$T_d(5)=((5+\sqrt{24})^d+(5-\sqrt{24})^d)/2$, a logarithmic degree suffices.
This bounds certificate degree, not the cost of obtaining the arithmetic
moments. It refines the earlier fixed Chebyshev filter; neither filtering nor
Christoffel minimization is claimed as a new general principle.

\paragraph{Certified moment intervals.}
If $m_j\in[\underline m_j,\overline m_j]$, select any rational normalized
polynomials $P_+,P_-$. For their upper polynomial $P_+^2$ and lower polynomial
$(x-c)P_-^2/(1-c)$, sum moment intervals with the sign of each coefficient.
The upper endpoint of the first sum bounds $H$ above; the lower endpoint of
the second bounds it below. Approximate optimization may choose the polynomials
but is not proof of their moments. Singular matrices need not be inverted,
and signed sums are not bounded by independently subtracting two upper bounds.

\section{Exact examples and finite verification}
At $p=2521,a=636,R=23$, the original measure is
\[
 3\delta_1+87\delta_{1/23}.
\]
The boundary module is three labelled copies of $\mathbb Q$ and 87 labelled
copies of $\mathbb Q[z]/z^2$: total dimension 177, socle dimension 90,
but $\operatorname{rank}(\epsilon s)=3$. For degree one,
$P(x)=(23x-1)/22$ gives $L_1=U_1=3$. The split $e=(2,1,1)$,
$f=(1,1,1)$ has five matching pairs, but its inverse-fibre weights give the
same measure and the same count three.

At the empty shell $p=241,a=64,R=15$, the measure is
$8\delta_{1/3}+12\delta_{1/5}+6\delta_{1/15}$ and the complete boundary socle
still has 26 labelled directions. Its augmentation comparison is zero.
The calculated bracket is
\[
 \boxed{-\frac{12}{425}\le H\le\frac{180}{463}<1,}
\]
which certifies $H=0$ in that shell, not failure at the prime.

The exhaustive comparison uses all 9531 first-half shells of the 54 primes
$p\equiv1\pmod{12}$ with $p\le1500$: 412838 original channel-labelled states,
1920 successful states in 685 shells, and 8846 empty shells. Degree-one
brackets determine the exact count in 9530 shells; degree two resolves the
remaining shell $p=1093,a=280,R=27$. There
\[
 [L_1,U_1]=[45982/7737,681/94],\qquad
 [L_2,U_2]=[6,1132982/188317].
\]
All 685 occupied shells are positively certified at degree at most one.
This is not new prime coverage or a universal lower bound. The implementation
computes the actual moments by finite enumeration; no faster algorithm is
claimed. Separate exact fixtures test singular cases in 80 small measures,
all 64 three-prime exponent tuples through three, coordinatewise reduction,
split weights, source secants, and six false-formula rejection tests.

\section{Remaining pointwise target and formalization scope}
For the full atlas of a prescribed prime, ES failure would mean that every
socle top monomial has a nonzero valuation degree, equivalently $\epsilon s=0$.
The new positive test seeks $L_d>0$ on some shell or finite combined atlas.
The proofs do not force that inequality uniformly. A separate arithmetic
bound on the original moments is still required; merely constructing the
boundary module or checking positive Gram matrices does not supply it.

The coefficient modules, their support labels, original exponent vectors,
channel maps, moment observation kernel, singular correction kernels and
prime-wise reductions remain explicit. The new finite constructions do not
inherit a Frobenius, a theta norm, or a zeta-zero interpretation by terminology.
No source publication or unmerged Zeta branch is altered. Source claims of
Lean/CI status are not a new Lean audit of this application.
